\documentclass[12pt,a4paper,oneside]{article}
\usepackage[brazil]{babel}
\usepackage[utf8]{inputenc}
\usepackage{olymp}
\usepackage{color}
\usepackage[dvipsnames]{xcolor}
\usepackage{listings}

\setcounter{problem}{4}

\begin{document}

\begin{problem}{Soma exata}{somar}{2}
 
Dada uma lista de $N$ números inteiros, seu programa deve determinar se existem dois elementos em posições distintas da lista cuja soma seja exatamente igual a um determinado valor $X$.

%\Notes
%
%Observações, se tiver.

\InputFile

A entrada contém três linhas: a primeira contém um valor inteiro $N$, indicando a quantidade de números na lista; a segunda contém os $N$ números da lista, que são números inteiros $V_i$ separados por espaço em branco; e a terceira linha contém um valor inteiro $X$.

Restrições: $1 \leq N \leq 1000$, $-10^6 \leq X \leq 10^6$, $-10^6 \leq V_i \leq 10^6$.



\OutputFile

Seu programa deve gerar apenas uma linha de saída, contendo uma única palavra, indicando se existem dois elementos em posições distinas da lista cuja soma seja exatamente igual a $X$. Se existirem, escreva ``SIM'', caso contrário, escreva ``NAO''.

%\Notes
%Nesta questão, seu programa deve usar a função \texttt{fatorial} dada %acima exatamente como está, não a modifique! Sua
%tarefa é apenas criar o programa com a função \texttt{main}.

\Example

\begin{example}
\exmp{
10
1 2 3 4 5 6 7 8 9 10
11
}{
SIM
}%
\end{example}

\begin{example}
\exmp{
10
1 2 3 4 5 6 7 8 9 10
20
}{
NAO
}%
\end{example}

\begin{example}
\exmp{
3
1 4 6
8
}{
NAO
}%
\end{example}

%Comentários sobre o exemplo, se necessário.

\end{problem}

\end{document}
